% Generated draft: requires mathematical and source review.
\documentclass[11pt]{article}
\usepackage[margin=25mm]{geometry}
\usepackage{fontspec}
\setmainfont{Latin Modern Roman}
\usepackage{amsmath,amssymb,mathtools,unicode-math}
\setmathfont{Latin Modern Math}
\usepackage{enumitem,needspace,xurl,hyperref}
\setlength{\parindent}{0pt}
\setlength{\parskip}{5pt}
\setlength{\emergencystretch}{3em}
\newcommand{\N}{\mathbb{N}}
\newcommand{\Z}{\mathbb{Z}}
\newcommand{\Q}{\mathbb{Q}}
\newcommand{\R}{\mathbb{R}}
\newcommand{\C}{\mathbb{C}}
\newcommand{\F}{\mathbb{F}}
\newcommand{\A}{\mathbb{A}}
\newcommand{\PP}{\mathbb{P}}
\newcommand{\E}{\mathbb{E}}
\newcommand{\eps}{\varepsilon}
\newcommand{\dd}{\,\mathrm{d}}
\begin{document}
\section*{UnsolvedMath ID: 6200016}
\textbf{Source problem code:} catalogue entry
\subsection{Definitions and mathematical statement}
For a nonempty finite-dimensional compact metric space \(Z\), let \(\dim Z\) denote covering dimension. A Boltyansky compactum is one satisfying
\[
\dim(Z\times Z)=2\dim Z-1.
\]
Dranishnikov's conjecture is that a compactum \(Z\) which is a group \(G\)'s \(Z\)-boundary never has this property. In the finite-dimensional setting its exact negation-of-Boltyansky target is
\[
\forall(G,Z)\text{ with }Z\text{ a }Z\text{-boundary of }G,
\qquad \dim(Z\times Z)\ne2\dim Z-1.
\]

The boundary notion refers to a \(Z\)-structure, not merely an arbitrary compactification remainder. In a common finite-dimensional convention it is a compactification \(\overline X=X\cup Z\) by a compact absolute retract, with \(Z\) a closed \(Z\)-set, a proper cocompact action of \(G\) on \(X\), and the nullity condition. The \(Z\)-set condition means there is a homotopy \(H\) of \(\overline X\) with \(H_0\) the identity and \(H_t(\overline X)\cap Z=\varnothing\) for \(t>0\). Nullity says that for each compact \(C\subseteq X\) and open cover of \(\overline X\), all but finitely many translates \(gC\) lie inside a member of the cover. The source does not specify its exact retract/dimension or action-freeness convention, so these choices need verification before fixing the quantified class.

The supplied source reports a proved special case when \(Z\) is a Markov compactum with all indicated building-block maps \(K_\sigma\to\sigma\) isomorphic. The inverse-limit construction and map-isomorphism convention are absent; this homogeneous-block qualification is retained rather than expanded to every possible meaning of Markov compactum. Reference [29] is not supplied. The research summary describes the general case as conjectural and explicitly says that a modern resolution was not verified.
\subsection{Short English statement}
A group \(Z\)-boundary should never have product dimension one below twice its own dimension. The source reports a homogeneous Markov-compactum special case, but leaves the precise boundary and building-block conventions to its missing definitions.
\subsection{Sources}
\begin{itemize}
\item Local UnsolvedMath registry, permanent ID \(6200016\), exact source code \texttt{catalogue entry}, ``Boundaries of Groups and Kleinian Groups --- Problem 16.'' Proposer Alexander Dranishnikov, 2005.
\item Kapovich, \emph{Problems on Boundaries of Groups and Kleinian Groups} (2005), Problem 16, recorded PDF page 6; \url{https://www.math.ucdavis.edu/~kapovich/EPR/problems.pdf}. Recorded access 2026-07-29; pdf-text-normalized extraction.
\item The source attributes the homogeneous Markov-compactum special case to reference [29], whose bibliography entry is not included. The OPEN-TRIAGE research summary supplies no complete later reference.
\end{itemize}
\end{document}
